\section{Model Mesin dan Kedudukan Bahasa Assembly}
\label{sec:model-mesin}

Sebelum menulis satu instruksi pun, mahasiswa perlu memiliki gambar
mental tentang apa yang sebenarnya terjadi di dalam komputer.
Bagian ini membangun gambar mental tersebut.

\subsection{Model Mesin von Neumann}

Hampir seluruh komputer modern mengikuti model yang dirumuskan John
von Neumann pada tahun 1945. Model ini memiliki empat komponen
utama.

\begin{enumerate}
    \item \term{Unit pemroses} (\eng{central processing unit}, CPU)
          yang menjalankan instruksi.
    \item \term{Memori} yang menyimpan baik instruksi maupun data.
    \item \term{Unit masukan} yang membawa data dari dunia luar.
    \item \term{Unit keluaran} yang menyampaikan hasil kembali ke
          dunia luar.
\end{enumerate}

Ciri pembeda model von Neumann adalah \term{prinsip
stored-program}: instruksi dan data disimpan bersama-sama di memori
yang sama, tanpa pemisahan jenis. Akibatnya, sebuah program dapat
memperlakukan instruksi sebagai data dan sebaliknya. Sifat inilah
yang membuat program dapat dimuat dari cakram ke memori lalu
dijalankan, dan sekaligus menjadi dasar teknik lanjutan seperti
\eng{self-modifying code} serta sebagian besar teknik eksploitasi
perangkat lunak.

\begin{figure}[H]
    \centering
    \begin{tikzpicture}[
        blok/.style={
            draw=obeBlue, line width=0.9pt, fill=obeBlueLight,
            rounded corners=3pt, minimum height=1.25cm,
            align=center, font=\small\bfseries, inner sep=6pt
        },
        panah/.style={-{Stealth[length=2.5mm]}, line width=1pt, obeTeal}
    ]
        \node[blok, minimum width=3.4cm] (cpu) at (0,0)
            {CPU\\{\normalfont\scriptsize unit pemroses}};
        \node[blok, minimum width=3.4cm, fill=obeAmberLight,
              draw=obeAmber] (mem) at (6.4,0)
            {Memori\\{\normalfont\scriptsize instruksi + data}};
        \node[blok, minimum width=2.8cm, fill=obeTealLight,
              draw=obeTeal] (in) at (-4.6,1.5)
            {Masukan};
        \node[blok, minimum width=2.8cm, fill=obeTealLight,
              draw=obeTeal] (out) at (-4.6,-1.5)
            {Keluaran};

        \draw[panah] (cpu) -- node[above, font=\scriptsize]
            {bus alamat} node[below, font=\scriptsize] {bus data}
            (mem);
        \draw[panah] (in) -- (cpu);
        \draw[panah] (cpu) -- (out);
    \end{tikzpicture}
    \caption{Model mesin von Neumann: CPU, memori bersama untuk
    instruksi dan data, serta unit masukan dan keluaran.}
    \label{fig:von-neumann}
\end{figure}

\subsection{Siklus Instruksi}

CPU bekerja dengan mengulang tiga langkah yang sama tanpa henti.

\begin{enumerate}
    \item \term{Fetch} (\eng{pengambilan}) --- CPU mengambil
          instruksi berikutnya dari memori. Alamat instruksi
          ditunjuk oleh register \reg{IP} pada 8086.
    \item \term{Decode} (\eng{penerjemahan}) --- CPU menentukan
          instruksi apa yang harus dikerjakan dan operand apa yang
          diperlukan.
    \item \term{Execute} (\eng{pelaksanaan}) --- CPU menjalankan
          operasinya, misalnya menjumlahkan dua bilangan atau
          memindahkan data.
\end{enumerate}

Pada 8086 terdapat langkah tambahan yaitu \eng{store}
(\eng{write-back}), ketika hasil eksekusi ditulis kembali ke register
atau memori. Siklus ini dibahas lebih rinci pada
Bab~\ref{chap:arsitektur-8086}.

\begin{infobox}
Perhatikan bahwa \emph{data} juga harus diambil dari memori, bukan
hanya instruksi. Karena itu satu instruksi seperti
\asm{ADD AX, [BX]} menimbulkan lebih dari satu akses memori: satu
untuk mengambil instruksi, satu atau lebih untuk mengambil data.
Jumlah akses memori inilah yang menjadi salah satu ukuran biaya
sebuah instruksi.
\end{infobox}

\subsection{Kedudukan Bahasa Assembly}

Bahasa assembly adalah representasi tekstual dari bahasa mesin.
Ketika Anda menulis:

\begin{lstlisting}[language=MASM, numbers=none]
MOV AX, 0005h
\end{lstlisting}

assembler menerjemahkannya menjadi deretan byte tertentu, misalnya
\hex{B8 05 00}. Ketika Anda menulis \asm{ADD AX, BX}, assembler
menghasilkan \hex{03 C3}. Setiap mnemonic dapat dipetakan ke satu
pola bit yang pasti.

\begin{table}[H]
    \centering
    \caption{Tiga baris assembly dan padanan kode mesinnya.}
    \label{tab:mnemonic-opcode}
    \small
    \begin{tabular}{@{}llp{5.6cm}@{}}
        \toprule
        \textbf{Assembly} & \textbf{Kode mesin} &
        \textbf{Keterangan} \\
        \midrule
        \asm{MOV AX, 0005h} & \hex{B8 05 00} &
        Nilai segera 16-bit disimpan \eng{little-endian}: byte rendah
        lebih dahulu \\
        \addlinespace
        \asm{ADD AX, BX} & \hex{03 C3} &
        Satu byte \eng{opcode} ditambah satu byte ModR/M \\
        \addlinespace
        \asm{INT 21h} & \hex{CD 21} &
        Instruksi interupsi dengan nomor vektor sebagai operand \\
        \bottomrule
    \end{tabular}
\end{table}

Karena hubungannya yang satu-ke-satu inilah bahasa assembly disebut
\term{mnemonic} untuk bahasa mesin: ia tidak menambah kemampuan
baru, hanya membuat kode mesin dapat dibaca manusia. Assembler
modern menyediakan bantuan seperti makro dan direktif, tetapi
kemampuan inti yang ada di dalamnya tetap sama.

\begin{infobox}
Urutan byte pada \hex{B8 05 00} mengikuti konvensi
\term{little-endian}: byte yang paling tidak bermakna disimpan pada
alamat terendah. Nilai \hex{0005} karena itu tersimpan sebagai
\hex{05} lalu \hex{00}. Konvensi ini berlaku untuk seluruh data
16-bit dan 32-bit pada keluarga x86 dan akan sering muncul pada
Bab~\ref{chap:segmentasi} serta Bab~\ref{chap:format-instruksi}.
\end{infobox}
